\begin{recipe}{Appelbollen}
\kicker[Dessert,Zoet]{Dessert · zoet}
\index[register]{Appelbollen}
\index[register]{Appel}
\index[register]{Bladerdeeg}
\index[register]{Kaneel}

\meta{45 minuten \dvd Voor 5 personen \dvd Oven 200\,°C}

\lettrine{M}{et} dit eenvoudige basisrecept maak je zelf appelbollen met bladerdeeg,
gevuld met kaneelsuiker. Handig als er nog een paar appels over zijn, en leuk om
samen te doen in het weekend.

\blockrule{Ingrediënten}
\begin{ingredients}
  \ing{5}{kleine, stevige appels (bijvoorbeeld Jonagold, Elstar of Goudreinet)}\index[register]{Appel}
  \ing{5}{vierkante plakjes bladerdeeg (ontdooid)}\index[register]{Bladerdeeg}
  \ing{2 el}{suiker}
  \ing{1 tl}{kaneel}\index[register]{Kaneel}
  \ing{1}{ei, losgeklopt (om te bestrijken)}
  \ingb{extra kristalsuiker, voor de buitenkant}
\end{ingredients}

\blockrule{Bereiding}
\tip{Serveer de appelbollen het liefst lauwwarm.}
\begin{steps}
  \item Verwarm de oven voor op 200\,°C en bekleed een bakplaat met bakpapier.
  \item Schil de appels en verwijder het klokhuis met een appelboor.
  \item Meng de suiker en kaneel in een schaaltje en rol de appels hier
        doorheen zodat ze rondom bedekt zijn. Meng eventueel wat extra suiker
        door het gat.
  \item Leg elk plakje bladerdeeg plat neer en plaats een appel in het
        midden. Vouw de punten van het deeg strak naar boven om de appel heen
        en druk de naden goed dicht met een klein drupje water, zodat er
        tijdens het bakken geen sap kan ontsnappen.
  \item Draai de bol om en leg hem met de gladde kant naar boven op de
        bakplaat. Bestrijk de bollen met het losgeklopte ei en strooi er nog
        wat extra suiker overheen.
  \item Bak de appelbollen 25 tot 30 minuten tot ze goudbruin en gaar zijn.
\end{steps}
\end{recipe}
